\section{The arithmetic Floquet normal form and critical recurrence law}\label{sec:floquet}
Put \[
 \alpha=1+\sqrt2,\quad \lambda=\alpha^2,\quad
 T=\begin{pmatrix}3&2\\4&3\end{pmatrix}.
\] Let K tend to infinity through odd primes, and define positive integers N,F by \[
 N+F\sqrt2=\alpha^K,\qquad N^2-2F^2=-1.
\] The moduli N may be composite. On the odd subspace of the standard N-dimensional torus quantization, let \(U_{N}\) be the phase-normalized quantization of T described in Subsection~\ref{sec:floquet}.1, and let \[
 B_N f(x)=2\cos(2\pi x/N)f(x),\quad
 U_{\varepsilon,N}=e^{i\varepsilon B_N}U_N,\quad
 d_N=(N-1)/2.
\]

\textbf{Theorem A (critical resonance law).} For every fixed real \(\tau\) and every integer m, with \(\varepsilon=\tau/\sqrt{2K}\), \[
 \lim_{K\longrightarrow\infty}
 \frac{1}{d_N}\operatorname{Tr}_{-}
       U_{\varepsilon,N}^{mK}
       =\exp(-m^2\tau^2/2).                                      \tag{FN.A}
\] Consequently the empirical eigenvalue measure of \(U_{\varepsilon,N}^K\) converges weakly on the unit circle to the wrapped normal law. For \(\tau\) nonzero its angular density on {[}-\(\pi\),\(\pi\)) is \[
 w_\tau(\theta)=\frac{1}{\sqrt{2\pi}|\tau|}
       \sum_{\ell\in\mathbb Z}
       \exp\!\left[-\frac{(\theta+2\pi\ell)^2}{2\tau^2}\right].   \tag{FN.B}
\] For \(\tau=0\) the limit is the point mass at 1. In particular the normalized resonance echo obeys \[
 d_N^{-2}\left|\operatorname{Tr}_{-}U_{\varepsilon,N}^{mK}\right|^2
       \longrightarrow e^{-m^2\tau^2}.                          \tag{FN.C}
\] This is a deterministic limit: there is no parameter average in (FN.A).

The principal finite estimate, proved before any limit is taken, is \[
 \left\|V_\varepsilon U_{\varepsilon,N}V_\varepsilon^*
            -e^{i\varepsilon\overline B}U_N\right\|_{2,-}
       \leq \frac{12\varepsilon^2}{\sqrt K},\qquad K\geq11,        \tag{FN.D}
\] for an explicit unitary \(V_{\varepsilon}\). Here the norm is normalized Hilbert--Schmidt norm on the entire odd subspace. It does not control an individual level or an operator-norm spectral gap. The finite comparison error in (FN.A) is at most \(6|m|\tau^2/\sqrt K\); an additional Gaussian characteristic-function convergence error is \(o(1)\), not assigned that rate.

\subsection{Arithmetic and the fixed classical observable}\label{floquet:arithmetic-and-the-fixed-classical-observable}

The recurrence \((N_j,F_j)=(N_{j-1}+2F_{j-1},N_{j-1}+F_{j-1})\), starting at \((1,0)\), shows that \(N\) is odd. In \(\mathbb Z/N\mathbb Z\), put \[
 s=F^{-1},\qquad s^2=2,\qquad h=3+2s,\qquad q=-F.
\] Since \((1+s)^K=N+Fs=1\), we have \(h^K=1\). Also \(h-1=2(1+s)\) is a unit. Modulo every prime \(\ell\mid N\), the order of \(h\) is precisely \(K\): it divides the prime \(K\) and is not one. Hence \(1-h^j\) is a unit for \(1\leq j<K\), and \(1+h^j\) is a unit for every \(j\), because \(K\) is odd. In particular, \(\ell\equiv1\pmod K\) and \(N\equiv1\pmod{2K}\). Neither squarefreeness nor primality of \(N\) is required.

Use the counting-measure inner product and centered Weyl operators \[
 W(a,b)f(x)=e_N(a(x-b/2))f(x-b),\qquad
 S_u f(x)=f(u^{-1}x),\qquad Rf(x)=f(-x).
\] Divisions in Weyl exponents mean inverses in \(\mathbb Z/N\mathbb Z\). The multiplication and covariance formulas are \[
 W(a,b)W(c,d)=e_N((ad-bc)/2)W(a+c,b+d),
\] \[
 S_uW(a,b)S_u^*=W(u^{-1}a,ub),\qquad RW(v)R=W(-v).
\] Individual Weyl translations need not preserve parity; their symmetrized sums do. On the odd subspace take \[
 U_0=S_h^{-1},\qquad B=W(q,1)+W(q,1)^*,\qquad
 U_\varepsilon=e^{i\varepsilon B}U_0.                            \tag{FN.1}
\] In these coordinates \(U_0^K=I\) exactly.

For completeness, this is unitarily equivalent to the fixed-map formulation of Theorem A, including the kick, not merely its compressed observable. Let \(C_tf(x)=e_N(tx^2/2)f(x)\), and let \(F_N\) have kernel \(e_N(xy)/\sqrt N\). Set \[
 M=\begin{pmatrix}-1/(2s)&1\\-1/2&-s\end{pmatrix},\qquad
 V=C_{s^{-1}}S_{1/2}F_N C_{2s}.
\] Its factors have Weyl covariance matrices respectively \[
 \begin{pmatrix}1&t\\0&1\end{pmatrix},\quad
 \begin{pmatrix}u^{-1}&0\\0&u\end{pmatrix},\quad
 \begin{pmatrix}0&1\\-1&0\end{pmatrix}.
\] Multiplication gives covariance \(M\), \(\det M=1\), and \(T=M\operatorname{diag}(h,h^{-1})M^{-1}\). Define \(U_N=VU_0V^*\). The index of \(V^*W(1,0)V\) is \(M^{-1}(1,0)=(-s,1/2)\). Further conjugation by \(S_2\), which commutes with \(U_0\), changes this to \((-s/2,1)=(-F,1)\pmod N\). Thus \(S_2V^*\) conjugates the entire physical kicked propagator to (FN.1). These unitaries all preserve parity.

With the usual semiclassical convention \(e^{i\,\mathrm{action}/\hbar}\), the classical kick amplitude associated with \(\varepsilon\) is proportional to \(\hbar\varepsilon\). Here \(\hbar\asymp1/N\). At the critical scale the amplitude is of order \(1/(N\sqrt K)\), not a fixed classical perturbation.

\subsection{Two summable phase estimates}\label{floquet:two-summable-phase-estimates}

Write \(C(A)=U_0 A U_0^*\), and \[
 B_j=C^j(B)=W(v_j)+W(-v_j),\qquad v_j=(qh^j,h^{-j}),
 \qquad \overline B=K^{-1}\sum_{j=0}^{K-1}B_j.                    \tag{FN.2}
\] The conjugation C shifts the label cyclically, so \(C(\bar B)=\bar B\). Define \[
 a_j=\left|e_N(q(h^j-h^{-j}))-1\right|,\qquad
 c_j=\frac{K-1-2j}{2K},\qquad G=\sum_{j=0}^{K-1}c_jB_j.          \tag{FN.3}
\] Then \(a_0=0,\ a_{K-j}=a_j\). Direct subtraction of the coefficients gives the exact cohomological equation \[
 G-C(G)=B-\overline B.                                         \tag{FN.4}
\] The coefficients have zero mean. Importantly, they are small at the half-period at which the large commutator phases accumulate.

Let \(P_{t}\) denote the second Pell sequence, \(\alpha^t=Q_t+P_t\sqrt2\). For the signed representative \(\delta\) of j mod K one has \[
 q(h^j-h^{-j})=-2P_{2\delta}\pmod N.                            \tag{FN.5}
\] Indeed \(h^\delta=(3+2s)^\delta=Q_{2\delta}+P_{2\delta}s\), and \(Fs=1\). For positive \(\delta=(K-1)/2-\ell\), \[
 \frac{P_{2\delta}}{N}
 =\frac{\lambda^{-\ell}}{\sqrt2\,\alpha}
       \frac{1-\lambda^{-2\delta}}{1-\lambda^{-K}}.              \tag{FN.6}
\] Put \(r=\lambda^{-1}\), \(A=4\pi/(\sqrt2\,\alpha)\). Using \(|e^{it}-1|\leq|t|\) and summing both signed representatives proves \[
 L_K:=\sum_{j=0}^{K-1}a_j
 \leq\frac{2A}{(1-r)(1-\lambda^{-K})}
 =\frac{2\sqrt2\pi}{1-\lambda^{-K}}<9\quad(K\geq3).              \tag{FN.7}
\] There is a sharper weighted estimate. At j=(K-1)/2-ell the weights at j and K-j have absolute values ell/K and (ell+1)/K. Therefore \[
 \sum_j |c_j|a_j
 \leq\frac{A}{K(1-\lambda^{-K})}\sum_{\ell\geq0}(2\ell+1)r^\ell
 =\frac{2\pi}{K(1-\lambda^{-K})}.                              \tag{FN.8}
\] The final identity follows from \(A(1+r)/(1-r)^2=2\pi\). The estimates are analytic inequalities; the regression's decimal phase values are only finite checks of them.

\subsection{Quadratic lifting and orthogonal commutator supports}\label{floquet:quadratic-lifting-and-orthogonal-commutator-supports}

\textbf{Lemma 1 (lifting).} Fix m. For all sufficiently large K, every signed m-term reciprocal relation \[
 \sum_{\nu=1}^m\epsilon_\nu h^{j_\nu}
 =\sum_{\nu=1}^m\epsilon_\nu h^{-j_\nu}=0\pmod N,
 \qquad \epsilon_\nu\in\{-1,1\},                               \tag{FN.9}
\] can be cyclically rotated to a polynomial \(P\) with \(P(\lambda)=0\). Its coefficient \(\ell^1\)-norm is at most \(m\), and its degree is at most \(K-\lceil K/m\rceil\). For \(m=4\), this holds for every odd prime \(K\geq11\), and the relation is then coefficientwise balanced.

\textbf{Proof.} Cancel opposite signs at identical labels. If anything remains, cut the cyclic support at a largest gap and rotate. There are at most \(m\) support points, giving the stated degree bound. Both congruences in (FN.9) survive. Divide \(P\) by \(Z^2-6Z+1\) over the integers. Its remainder \(aZ+b\) vanishes at \(h\) and \(h^{-1}\); their difference \(4s\) is a unit. Thus \(N\mid a\) and \(N\mid b\). A nonzero quadratic norm of \(P(\lambda)\) would be an integer divisible by \(N^2\), but \[
 |P(\lambda)P(\lambda^{-1})|
 \leq m^2\lambda^{K-\lceil K/m\rceil}<N^2                       \tag{FN.10}
